% ============================================================================
%  CLIPP — MILP Formulation  (copy-paste into Overleaf)
% ============================================================================
%  Requires:  amsmath, amssymb, amsthm (loaded by most templates).
%  Standalone compilation:  pdflatex clipp_formulation.tex
% ============================================================================
\documentclass[11pt]{article}
\usepackage[margin=2.5cm]{geometry}
\usepackage{amsmath,amssymb}
\usepackage{enumitem}
\usepackage{booktabs}
\usepackage{xcolor}
\usepackage{listings}

% ---- Python code style -----------------------------------------------------
\definecolor{codebg}{HTML}{F5F5F5}
\definecolor{codegreen}{HTML}{2E7D32}
\definecolor{codegray}{HTML}{757575}
\definecolor{codepurple}{HTML}{7B1FA2}
\definecolor{codeblue}{HTML}{1565C0}

\lstdefinestyle{python}{
    language=Python,
    backgroundcolor=\color{codebg},
    basicstyle=\ttfamily\footnotesize,
    keywordstyle=\color{codeblue}\bfseries,
    stringstyle=\color{codegreen},
    commentstyle=\color{codegray}\itshape,
    emphstyle=\color{codepurple},
    emph={model,solver,cp_model,CpModel,CpSolver,IntVar,BoolVar},
    numberstyle=\tiny\color{codegray},
    numbers=left,
    numbersep=6pt,
    frame=single,
    rulecolor=\color{codegray!40},
    breaklines=true,
    breakatwhitespace=false,
    tabsize=4,
    showstringspaces=false,
    captionpos=b,
    aboveskip=8pt,
    belowskip=8pt,
}
\lstset{style=python}

\newcommand{\setN}{\mathcal{N}}   % set of nodes
\newcommand{\setA}{\mathcal{A}}   % set of arcs
\newcommand{\setJ}{\mathcal{J}}   % set of streets
\newcommand{\setK}{\mathcal{K}}   % set of vehicles
\newcommand{\Jm}{\mathcal{J}_M}   % mandatory streets
\newcommand{\Jo}{\mathcal{J}_O}   % optional streets
\newcommand{\Jc}{\mathcal{J}_C}   % connector streets
\newcommand{\Aj}{A_j}             % arcs of street j
\newcommand{\depot}{S}

\title{CLIPP --- Mixed-Integer Linear Programming Formulation}
\author{}
\date{}

\begin{document}
\maketitle

%% -------------------------------------------------------------------
\section{Problem overview}
%% -------------------------------------------------------------------

The \textbf{Climate-Integrated Street Cleaning Routing Problem (CLIPP)}
plans routes for a heterogeneous fleet of cleaning vehicles on a directed
street network.  The objective is to maximise the total length of streets
cleaned while minimising water waste, subject to a shared time budget and
mandatory-coverage requirements.

%% -------------------------------------------------------------------
\section{Notation}
%% -------------------------------------------------------------------

\subsection{Sets}

\begin{table}[h]
\centering
\begin{tabular}{@{} l l @{}}
\toprule
Symbol & Description \\
\midrule
$\setN = \{0, \dots, N-1\}$       & Junctions (nodes) \\
$\setJ = \{0, \dots, M-1\}$       & Streets (edges) \\
$\Jm \subset \setJ$               & Mandatory streets \\
$\Jo \subset \setJ$               & Optional streets \\
$\Jc \subset \setJ$               & Connector streets ($\setJ = \Jm \cup \Jo \cup \Jc$) \\
$\setA$                            & Directed arcs (a two-way street generates two arcs) \\
$\Aj \subseteq \setA$             & Set of arcs corresponding to street $j$ \\
$\setK = \{0, \dots, C-1\}$       & Vehicles \\
$\delta^+(u), \delta^-(u)$         & Outgoing / incoming arcs at node $u$ \\
\bottomrule
\end{tabular}
\end{table}

\subsection{Parameters}

\begin{table}[h]
\centering
\begin{tabular}{@{} l l @{}}
\toprule
Symbol & Description \\
\midrule
$T$            & Shared time limit for every vehicle (seconds) \\
$\depot$       & Depot junction \\
$\alpha$       & Coverage/efficiency trade-off weight, $0 \le \alpha \le 1$ \\
$L_{\max}$     & Total cleanable street length (mandatory + optional) \\
$W_{\max}$     & Worst-case waste if all cleanable streets use a Large vehicle \\
$c_a$          & Traversal time of arc $a$ (seconds) \\
$\ell_j$       & Length of street $j$ (metres) \\
$r_j$          & Cleaning requirement of street $j$ (litres\,/\,km) \\
$q_k$          & Water capacity of vehicle $k$ (litres\,/\,km) \\
$U_a$          & Traversal upper bound for arc $a$: $U_a = \min\!\bigl(\lfloor T / c_a \rfloor,\; 50\bigr)$ \\
$\mathcal{M}$  & Big-M constant: $\mathcal{M} = 1 + \sum_{a \in \setA} U_a$ \\
\bottomrule
\end{tabular}
\end{table}

%% -------------------------------------------------------------------
\section{Decision variables}
%% -------------------------------------------------------------------

\begin{align}
f_{a,k} &\in \mathbb{Z}_{\ge 0}
    &&\forall\, a \in \setA,\; k \in \setK
    \tag{flow variable}\\[4pt]
x_{j,k} &\in \{0,1\}
    &&\forall\, j \in \setJ \setminus \Jc,\; k \in \setK
    \tag{cleaning assignment}\\[4pt]
y_j     &\in \{0,1\}
    &&\forall\, j \in \setJ \setminus \Jc
    \tag{cleaned indicator}
\end{align}

\begin{itemize}[nosep]
  \item $f_{a,k}$ counts how many times vehicle~$k$ traverses arc~$a$.
  \item $x_{j,k} = 1$ iff vehicle~$k$ cleans street~$j$.
  \item $y_j = 1$ iff street~$j$ is cleaned by \emph{any} vehicle.
\end{itemize}

\medskip\noindent
\textbf{Note.}
If vehicle~$k$ cannot clean street~$j$ (i.e.\ $q_k < r_j$, or $j \in \Jc$),
$x_{j,k}$ is fixed to~$0$ and not created as a decision variable.

\paragraph{Implementation (OR-Tools CP-SAT).}
\begin{lstlisting}[caption={Decision variable creation (\texttt{clipp/model.py})}]
# f[a, k]: number of times vehicle k traverses arc a
def _create_flow_vars(model, A, C, arc_max):
    f = {}
    for k in range(C):
        for a in range(A):
            f[a, k] = model.new_int_var(0, arc_max[a], f"f_{a}_{k}")
    return f

# clean[j, k]: BoolVar if vehicle k can clean street j, else fixed 0
def _create_clean_vars(model, M, C, streets, vehicles):
    clean = {}
    for j in range(M):
        s = streets[j]
        for k in range(C):
            if s["cat"] != "C" and vehicles[k]["cap"] >= s["req"]:
                clean[j, k] = model.new_bool_var(f"clean_{j}_{k}")
            else:
                clean[j, k] = 0  # cannot clean -> not a variable
    return clean

# cleaned[j]: BoolVar for non-connector streets
def _create_cleaned_vars(model, M, streets):
    cleaned = {}
    for j in range(M):
        if streets[j]["cat"] != "C":
            cleaned[j] = model.new_bool_var(f"cleaned_{j}")
        else:
            cleaned[j] = 0
    return cleaned
\end{lstlisting}

%% -------------------------------------------------------------------
\section{Constraints}
%% -------------------------------------------------------------------

% ---- (1) Flow conservation ----
\subsection{Flow conservation}

\begin{equation}\label{eq:flow}
\sum_{a \in \delta^+(u)} f_{a,k}
\;=\;
\sum_{a \in \delta^-(u)} f_{a,k}
\qquad \forall\, u \in \setN,\; k \in \setK
\end{equation}

Each vehicle's route forms a \emph{closed walk}: the total flow leaving
every junction equals the total flow entering it.  This guarantees that the
vehicle can always return to its starting point (the depot) without
``teleporting''.

\begin{lstlisting}[caption={Constraint~(\ref{eq:flow}): Flow conservation}]
def _add_flow_conservation(model, f, N, C, out_arcs, in_arcs):
    for k in range(C):
        for u in range(N):
            out_f = [f[a, k] for a in out_arcs.get(u, [])]
            in_f  = [f[a, k] for a in in_arcs.get(u, [])]
            if out_f or in_f:
                model.add(sum(out_f) == sum(in_f))
\end{lstlisting}

% ---- (2) Time limit ----
\subsection{Time limit}

\begin{equation}\label{eq:time}
\sum_{a \in \setA} c_a \, f_{a,k} \;\le\; T
\qquad \forall\, k \in \setK
\end{equation}

The total traversal time of every vehicle must not exceed the shared time
budget~$T$.

\begin{lstlisting}[caption={Constraint~(\ref{eq:time}): Time limit}]
def _add_time_constraints(model, f, A, C, T, arc_cost):
    for k in range(C):
        model.add(sum(arc_cost[a] * f[a, k] for a in range(A)) <= T)
\end{lstlisting}

% ---- (3)+(4) Traversal requirement ----
\subsection{Traversal requirement for cleaning}

\begin{equation}\label{eq:traverse}
\sum_{a \in \Aj} f_{a,k} \;\ge\; x_{j,k}
\qquad \forall\, j \in \setJ \setminus \Jc,\; k \in \setK
\end{equation}

If vehicle~$k$ cleans street~$j$ ($x_{j,k}=1$), it must have traversed at
least one arc of that street.  A vehicle cannot clean a street it has not
visited.

\begin{lstlisting}[caption={Constraint~(\ref{eq:traverse}): Traversal requirement}]
def _add_traversal_requirements(model, f, clean, M, C, streets, street_arcs):
    for j in range(M):
        if streets[j]["cat"] == "C":
            continue
        for k in range(C):
            if not isinstance(clean[j, k], int):  # skip fixed-zero vars
                traversals = [f[a, k] for a in street_arcs[j]]
                if traversals:
                    model.add(sum(traversals) >= clean[j, k])
\end{lstlisting}

% ---- (5) Cleaned linkage ----
\subsection{Cleaned-indicator linkage}

\begin{align}
x_{j,k} &\;\Rightarrow\; y_j
    \qquad &&\forall\, j \in \setJ \setminus \Jc,\; k \in \setK
    \label{eq:link-impl}\\[4pt]
y_j &\;\Rightarrow\; \bigvee_{k \in \setK} x_{j,k}
    \qquad &&\forall\, j \in \setJ \setminus \Jc
    \label{eq:link-or}
\end{align}

Together these encode the logical equivalence
$y_j \Leftrightarrow \bigvee_k x_{j,k}$.
Constraint~\eqref{eq:link-impl} ensures that the cleaned indicator is
activated whenever any vehicle cleans the street.
Constraint~\eqref{eq:link-or} prevents $y_j = 1$ unless at least one
vehicle actually cleans street~$j$.

\begin{lstlisting}[caption={Constraints~(\ref{eq:link-impl})\,--\,(\ref{eq:link-or}): Cleaned linkage}]
def _add_cleaned_linkage(model, clean, cleaned, M, C, streets):
    for j in range(M):
        if streets[j]["cat"] == "C":
            continue
        bool_vars = [
            clean[j, k] for k in range(C)
            if isinstance(clean[j, k], cp_model.IntVar)
        ]
        if bool_vars:
            for bv in bool_vars:
                model.add_implication(bv, cleaned[j])       # x_{j,k} => y_j
            model.add_bool_or(bool_vars + [cleaned[j].Not()])# y_j => OR_k x_{j,k}
        else:
            model.add(cleaned[j] == 0)
\end{lstlisting}

% ---- (6) Single cleaner ----
\subsection{Single-vehicle cleaning}

\begin{equation}\label{eq:single}
\sum_{k \in \setK} x_{j,k} \;\le\; 1
\qquad \forall\, j \in \setJ \setminus \Jc
\end{equation}

Each street is cleaned by \emph{at most one} vehicle.  While the problem
statement allows redundant cleaning, the solver enforces single-vehicle
cleaning to avoid unnecessary water waste (the street's length is scored
only once regardless).

\begin{lstlisting}[caption={Constraint~(\ref{eq:single}): Single-vehicle cleaning}]
def _add_single_cleaner(model, clean, M, C, streets):
    for j in range(M):
        if streets[j]["cat"] == "C":
            continue
        bool_vars = [
            clean[j, k] for k in range(C)
            if isinstance(clean[j, k], cp_model.IntVar)
        ]
        if len(bool_vars) > 1:
            model.add(sum(bool_vars) <= 1)
\end{lstlisting}

% ---- (7) Mandatory coverage ----
\subsection{Mandatory coverage}

\begin{equation}\label{eq:mandatory}
y_j = 1
\qquad \forall\, j \in \Jm
\end{equation}

Every mandatory street must be cleaned.  A solution that violates this
constraint is infeasible.

\medskip\noindent
\textbf{Relaxation.}
If a mandatory street is provably unreachable within the time budget (detected
via a Dijkstra-based pre-check), its hard constraint is relaxed to a soft
one --- it is still scored if cleaned, but the solver will not be forced into
infeasibility.

\begin{lstlisting}[caption={Constraint~(\ref{eq:mandatory}): Mandatory coverage}]
def _add_mandatory_constraints(model, cleaned, M, streets, relaxed_mandatory):
    for j in range(M):
        if streets[j]["cat"] == "M":
            if j in relaxed_mandatory:
                continue  # provably unreachable -- keep as soft
            assert isinstance(cleaned[j], cp_model.IntVar)
            model.add(cleaned[j] == 1)
\end{lstlisting}

% ---- (8) Depot inclusion ----
\subsection{Depot inclusion (Big-M)}

\begin{equation}\label{eq:depot}
\sum_{a \in \setA} f_{a,k}
\;\le\;
\mathcal{M} \cdot \sum_{a \in \delta^+(\depot)} f_{a,k}
\qquad \forall\, k \in \setK
\end{equation}

If vehicle~$k$ uses any arc at all (left-hand side $> 0$), then it must
have at least one outgoing traversal from the depot.  Combined with flow
conservation~\eqref{eq:flow}, this forces the vehicle's closed walk to
pass through the depot.  When the depot has zero outgoing flow, the
right-hand side becomes~$0$, prohibiting any arc usage --- i.e.\ an
unused vehicle simply stays at the depot.

\begin{lstlisting}[caption={Constraint~(\ref{eq:depot}): Depot inclusion (Big-M)}]
def _add_depot_inclusion(model, f, A, C, S, out_arcs, arc_max):
    big_M_val = sum(arc_max) + 1
    for k in range(C):
        depot_out_arcs = out_arcs.get(S, [])
        if depot_out_arcs:
            depot_out_sum  = sum(f[a, k] for a in depot_out_arcs)
            total_flow_sum = sum(f[a, k] for a in range(A))
            model.add(total_flow_sum <= big_M_val * depot_out_sum)
\end{lstlisting}

%% -------------------------------------------------------------------
\section{Objective function}
%% -------------------------------------------------------------------

\begin{equation}\label{eq:obj}
\max \quad
\alpha \cdot
\underbrace{%
    \frac{\sum_{j \in \setJ \setminus \Jc} \ell_j \, y_j}{L_{\max}}
}_{\text{Coverage}}
\;+\;
(1-\alpha) \cdot
\underbrace{%
    \left(
    1 -
    \frac{\sum_{j \in \setJ \setminus \Jc}
    \sum_{k \in \setK}
    \frac{(q_k - r_j) \, \ell_j}{1000}
    \, x_{j,k}}{W_{\max}}
    \right)
}_{\text{Efficiency}}
\end{equation}

\begin{itemize}[nosep]
  \item Coverage is normalized by $L_{\max}$ (total cleanable length).
  \item Efficiency is normalized by $W_{\max}$ (worst-case waste under Large-vehicle cleaning).
  \item The parameter $\alpha$ controls the trade-off. When $\alpha=1$, only
        coverage matters; when $\alpha=0$, only efficiency matters.
\end{itemize}

\medskip\noindent
\textbf{Implementation detail.}
To keep the solver in integer arithmetic, objective coefficients are scaled
by a large constant; this does not change the optimal solution.

\begin{lstlisting}[caption={Objective function implementation}]
_OBJ_SCALE = 1_000_000_000

def _set_objective(model, clean, cleaned, M, C, streets, vehicles, alpha):
    L_max = sum(s["length"] for s in streets if s["cat"] != "C")
    W_max = sum(
        max(0, 30 - s["req"]) * s["length"] / 1000.0
        for s in streets if s["cat"] != "C"
    )
    cov_weight = (alpha / L_max) if L_max > 0 else 0.0
    eff_penalty_weight = ((1.0 - alpha) / W_max) if W_max > 0 else 0.0
    obj_terms = []

    for j in range(M):
        s = streets[j]
        if s["cat"] == "C":
            continue

        # Positive term: normalized coverage contribution
        if isinstance(cleaned[j], cp_model.IntVar):
            coeff = int(round(_OBJ_SCALE * cov_weight * s["length"]))
            if coeff > 0:
                obj_terms.append(coeff * cleaned[j])

        # Negative term: normalized efficiency loss
        for k in range(C):
            if not isinstance(clean[j, k], cp_model.IntVar):
                continue
            excess = vehicles[k]["cap"] - s["req"]
            if excess > 0 and eff_penalty_weight > 0:
                waste_liters = excess * s["length"] / 1000.0
                penalty = int(round(_OBJ_SCALE * eff_penalty_weight * waste_liters))
                if penalty > 0:
                    obj_terms.append(-penalty * clean[j, k])

    model.maximize(sum(obj_terms) if obj_terms else 0)
\end{lstlisting}

%% -------------------------------------------------------------------
\section{Subtour elimination (iterative)}
%% -------------------------------------------------------------------

Flow conservation alone does not prevent \emph{subtours} — disconnected
cycles that never visit the depot.  These are eliminated lazily via an
iterative callback:

\begin{enumerate}[nosep]
  \item Solve the model.
  \item For each vehicle~$k$, compute the weakly-connected components of its
        arc-flow multigraph.
  \item For every component~$\mathcal{C}$ that does \emph{not} contain the
        depot, add the cut:
        \begin{equation}\label{eq:subtour}
        \sum_{\substack{a \in \setA :\\[1pt] \mathrm{head}(a) \in \mathcal{C} \\[1pt] \mathrm{tail}(a) \in \mathcal{C}}}
             f_{a,k}
        \;\le\;
        \biggl(\,1 + \!\!\sum_{\substack{a \in \setA :\\[1pt] \mathrm{head}(a) \in \mathcal{C} \\[1pt] \mathrm{tail}(a) \in \mathcal{C}}}\!\! U_a \biggr)
        \cdot
        \sum_{\substack{a \in \setA :\\[1pt] \mathrm{tail}(a) \notin \mathcal{C} \\[1pt] \mathrm{head}(a) \in \mathcal{C}}}
             f_{a,k}
        \end{equation}
        If no arcs enter the component, all internal arcs are set to~$0$.
  \item Re-solve with the new constraints (warm-started).
  \item Repeat until no subtours remain or the time budget is exhausted.
\end{enumerate}

\begin{lstlisting}[caption={Subtour elimination implementation}]
def _eliminate_subtours(model, f, f_vals, arcs, arc_max, A, C, S):
    subtours_found = False

    for k in range(C):
        total_used = sum(f_vals[a, k] for a in range(A))
        if total_used == 0:
            continue

        adj_k = build_adj_multigraph(k, f_vals, arcs)
        comps = connected_components(adj_k)

        for comp in comps:
            if S in comp:
                continue  # connected to depot -- OK

            internal = [
                a for a in range(A)
                if arcs[a][0] in comp and arcs[a][1] in comp
            ]
            entering = [
                a for a in range(A)
                if arcs[a][0] not in comp and arcs[a][1] in comp
            ]
            if not internal:
                continue

            if entering:
                big_M_sub = sum(arc_max[a] for a in internal) + 1
                model.add(
                    sum(f[a, k] for a in internal)
                    <= big_M_sub * sum(f[a, k] for a in entering)
                )
            else:
                for a in internal:
                    model.add(f[a, k] == 0)

            subtours_found = True
    return subtours_found
\end{lstlisting}

%% -------------------------------------------------------------------
\section{Post-processing: route reconstruction}
%% -------------------------------------------------------------------

After solving, each vehicle's arc-flow solution is converted into an
explicit junction sequence via \textbf{Hierholzer's algorithm} for Eulerian
circuits.  If any connected component of a vehicle's multigraph is
disconnected from the depot, a shortest-path bridge (depot $\to$ component
$\to$ depot) is inserted when the remaining time budget permits; otherwise
the component's arcs are dropped.

\end{document}
